\subsection{Key factors and practical guidance}\label{sec:factors}
\owner{Ky}

% Rubric: "discuss the key factors that affect the performance".

% [Ky]: ranked by how much each factor changed the outcome; every number is quoted from the section cited.
\paragraph{Key factors.} Ordered by the size of their effect in our experiments:
\begin{enumerate}\itemsep1pt
\item \textbf{Cluster shape against the method's assumption.} On D1 K-Means++ reaches ARI 0.481 and
DBSCAN 0.98; on D2 K-means joins belts thousands of kilometres apart (Section~\ref{sec:shape}). No
parameter setting repairs a wrong assumption about shape.
\item \textbf{Gaps between groups.} Density methods need empty space between clusters. Where there is
none they return one dominant cluster: 97.1\% of sales on D8, 96.1\% of pickups on D3 and 97\% of proteins
on D4 at the settings chosen without labels (Sections~\ref{sec:nogaps} and \ref{sec:density}).
\item \textbf{Dimensionality.} On D4 the ranking of DBSCAN and K-Means++ reverses between 3 and 4
principal components (Ablation~D); on D5, with 50 dimensions, DBSCAN reaches ARI 0.03 and HDBSCAN
calls 71\% of points noise (Section~\ref{sec:dims}).
\item \textbf{Uneven density.} One DBSCAN radius clusters 90\% of the dense south-west Pacific but 20\%
of the Americas on D2; HDBSCAN raises the Americas to 62\% (Section~\ref{sec:density}).
\item \textbf{Parameter choice without labels.} On D4 the rule's DBSCAN setting scores ARI 0.039 against
0.489 at the best setting; on D3 the knee and the silhouette-best radius differ threefold and give 13 and
2 clusters (Table~\ref{tab:cmp_d3d4}).
\item \textbf{Feature scaling.} Standardising D4 raises K-Means++ ARI from 0.497 to 0.731; on D8, raw
units turn K-means into price bands (Ablation~C).
\item \textbf{Values recorded in steps.} On D8, 107 of 124 HDBSCAN clusters are single storey bands on
the recorded data and none on jittered data (Ablation~E).
\item \textbf{Initialisation.} Over 30 seeds K-Means++ raises mean ARI on D4 from 0.485 to 0.561, and on
D2 it cuts the mean gap to the best inertia from 10.9\% to 4.8\%; on D3 and D8 it makes no difference to
the result (Ablation~A).
\item \textbf{Runtime.} It only matters at the size of D8, where HDBSCAN takes 212\,s and DBSCAN runs out
of memory (Section~\ref{sec:runtime}).
\end{enumerate}
% [Wei Yew]: evidence for the two guidance items handed to Ky.
% Fake-data check: D2 K=2 silhouette 0.544 vs 0.529 on shuffled coordinates (kmeans_k_selection.csv);
% D8 density methods below their null (comparison.csv). Recorded steps: D8 107 of 124 HDBSCAN clusters
% are single storey bands on recorded data, 0 of 23 on jittered data (hdbscan_grid.csv).
\paragraph{Practical guidance.}
\begin{enumerate}\itemsep1pt
\item \textbf{Choose by the structure of the data, not by a score.} Plot the data or a projection first.
If the groups are curved or long, no choice of $k$ will help K-means; if they are compact and touching,
or the data have more than a handful of dimensions, no choice of $\varepsilon$ will help DBSCAN.
\item \textbf{Do not pick settings by the silhouette alone.} It favours few, round clusters and settings
that call hard points noise. First restrict the search to what the analysis needs (a noise cap, a
minimum number of clusters), then compare within that range.
\item \textbf{Report the noise fraction and the largest cluster next to every silhouette.} A silhouette
computed after discarding 37\% of the points, or with 97\% in one cluster, says little.
\item \textbf{Compare every result against fake data.} Rerun the chosen setting on shuffled columns or
uniform points. It catches results with no real structure (D4 density methods, D8), though not results
with the wrong structure (D2).
\item \textbf{Check for values recorded in steps.} Tied values such as storey bands form dense spots that
density methods report as clusters; add jitter within each step and see whether the clusters remain.
\item \textbf{Standardise before K-means, and use K-means++ with several starts.} One K-means++ start is
not enough when the objective has many local optima (D2).
\item \textbf{Prefer a parameter whose unit means something, and measure distance correctly.} A radius
in kilometres can be checked by someone who knows the area; a $k$ cannot. Use distance on the sphere for
global coordinates (D2); a local projection is enough within a city (D3).
\end{enumerate}
